% GA-SUM-02 / EXP-1 — Notas del integrante 2 (Steffanny Loor)
% Guion oral aproximado: 4 minutos.
% Citas selectivas a fuentes técnicas; bibliografía compartida en referencias/loor.bib.

\section*{Guion de exposición — Steffanny Loor (4 minutos): comunicación y procesamiento paralelo}

\subsection*{Diapositiva 1. CAP: consistencia, disponibilidad y particiones (30 s)}
CAP nos ayuda a entender qué ocurre cuando se interrumpe la comunicación entre nodos.
Consistencia significa que las lecturas reflejan el estado más reciente;
la disponibilidad, que las solicitudes reciben una respuesta válida;
y tolerancia a particiones, que debemos considerar fallos de red.
 En una quorum queue de RabbitMQ con tres miembros se necesita una mayoría de dos.
 Si uno queda aislado, los otros dos pueden continuar, pero el miembro aislado no puede confirmar publicaciones p
 or su cuenta \cite{ahmed2023,myrhorodskyi2025,rabbitmq2026quorum}.

\subsection*{Diapositiva 2. PACELC (23 s)}
PACELC amplía este análisis.
Durante una partición, debemos considerar el compromiso entre
consistencia y disponibilidad. Incluso cuando la red funciona,
coordinar réplicas puede hacer que aumente la latencia.
En RabbitMQ, las quorum queues utilizan consenso y necesitan una mayoría.
 Por eso, no sería correcto clasificar a todo RabbitMQ de una sola manera:
 depende del tipo de cola y su configuración \cite{ahmed2023,myrhorodskyi2025,rabbitmq2026quorum}.

\subsection*{Diapositiva 3. Diagrama: dos formas de procesar un pedido (25 s)}
Este diagrama compara dos situaciones.
 En la primera, la aplicación espera que termine una
 operación antes de continuar. En la segunda, publica el
 pedido en una cola de RabbitMQ y puede seguir atendiendo s
 solicitudes. Varios consumidores comparten la cola y
 procesan pedidos diferentes. Esto permite distribuir el
 trabajo: si varios consumidores se ejecutan simultáneamente,
 los pedidos pueden procesarse en paralelo \cite{maharjan2023,derlatka2023,reinders2023}.

\subsection*{Diapositiva 4. Comunicación síncrona y asíncrona (12 s)}
La comunicación síncrona espera una respuesta para continuar;
la asíncrona permite avanzar sin esperar el procesamiento
del mensaje \cite{kazanavicius2023}.

\subsection*{Diapositiva 5. Comunicación síncrona (18 s)}
Una tienda necesita conocer si el pago fue aprobado antes de
confirmar la compra. La comunicación síncrona permite tomar
esa decisión, aunque puede introducir demoras \cite{kazanavicius2023}.

\subsection*{Diapositiva 6. Comunicación asíncrona (18 s)}
En la comunicación asíncrona, el emisor envía un mensaje y
no tiene que esperar a que el receptor termine de procesarlo.
Las colas almacenan mensajes pendientes y permiten que los componentes
trabajen en momentos diferentes.
Eso reduce su dependencia temporal, aunque requiere gestionar
los mensajes \cite{maharjan2023,derlatka2023}.

\subsection*{Diapositiva 7. Comparación (15 s)}
La diferencia principal es la espera. La comunicación síncrona
conviene cuando necesitamos una respuesta para decidir qué hacer.
La asíncrona conviene cuando una tarea puede quedar pendiente mientras
la aplicación continúa. Ninguna es mejor en todos los escenarios,
además, ser asíncrono no significa automáticamente ejecutar tareas en paralelo.

\subsection*{Diapositiva 8. Hilos, concurrencia y paralelismo (22 s)}
Un hilo es una unidad de ejecución dentro de un proceso. Hay
concurrencia cuando varias tareas progresan durante un mismo
período, aunque se alternen. Hay paralelismo cuando realmente
se ejecutan al mismo tiempo. Pensemos en tres trabajadores:
si se turnan, hay concurrencia; si trabajan simultáneamente,
hay paralelismo. En RabbitMQ, varios consumidores pueden recibir pedidos distintos;
solo existe paralelismo si los procesan al mismo tiempo. Tener varios consumidores no
lo garantiza por sí solo \cite{reinders2023}.

\subsection*{Diapositiva 9. Modelos de comunicación (20 s)}
En cliente-servidor, un proceso solicita un servicio y otro
lo atiende. En mensajería, los procesos intercambian mensajes,
a veces mediante un intermediario. En RPC, un proceso solicita
ejecutar una operación en otro. RabbitMQ funciona como intermediario
entre productores y consumidores y también permite construir patrones de petición
y respuesta \cite{zhang2023,kazanavicius2023}.

\subsection*{Diapositiva 10. Procesamiento asíncrono con RabbitMQ (18 s)}
En nuestro ejemplo, un cliente realiza un pedido,
la aplicación publica un mensaje y RabbitMQ lo dirige a una
cola. Después, uno de los consumidores recibe y procesa ese pedido.
Otros consumidores pueden atender pedidos distintos. Así, la aplicación
no espera a que cada trabajador termine. Si hay recursos para ejecutar
varios consumidores simultáneamente, diferentes pedidos pueden procesarse en
paralelo \cite{maharjan2023,derlatka2023}.

\subsection*{Diapositiva 11. Aplicación al PFC: SREI (27 s)}
En mi de fin de curso, SREI, en mi grupo tenemos dos servicios independientes
con Spring Boot. El servicio de eventos proporciona información y el
servicio de consulta la solicita mediante REST.
Elegimos comunicación síncrona para esta interacción porque el
consumidor necesita recibir los datos antes de responder a la consulta.
A diferencia del ejemplo de RabbitMQ, aquí no estamos distribuyendo trabajos
pendientes entre consumidores, sino obteniendo una respuesta directa.

\textbf{Fuentes utilizadas:} \cite{ahmed2023,myrhorodskyi2025,rabbitmq2026quorum,kazanavicius2023,maharjan2023,derlatka2023,reinders2023,zhang2023}.